{
    \setlength{\parindent}{0em}

    \par {\zihao{4}\bfseries 一、题目：\uline{\hfill \Title \hfill}}

    \par {\zihao{4}\bfseries 二、指导教师对开题报告的具体要求：}

    \setlength{\parindent}{2em}

    \par 本课题聚焦基于智能体的拖延症干预系统，针对知识工作者“启动困难”与“执行断裂”的核心痛点，探索对话式任务管理的新范式，兼具交互设计理论价值与数字健康管理应用前景。为保障毕业设计的顺利推进与最终质量，提出以下具体要求：

    \par （1）在国内外设计调研与理论建构方面，需紧密围绕对话式智能体系统的发展趋势，深入梳理人机协作范式、上下文感知计算及渐进式披露等相关理论。要求详实对比现有任务管理工具与 AIGC 应用在意图识别精度、多轮状态维持及情感化陪伴方面的典型案例，充分剖析当前系统在任务粒度控制与认知负荷平衡方面存在的结构性痛点，为后续智能体工作流设计确立坚实的理论支撑。

    \par （2）在设计概念与创新点方面，必须立足于移动端碎片化场景，严守低干扰与防挫败的设计底线。鼓励跳出传统表单化界面的经验主义框架，推导并建立一套严谨的“用户意图-任务结构”映射机制。要求结合智能体记忆增强与工具调用能力，将对话流编排与可视化进度管理有效融合，力求在人机协作式任务构建与长期记忆持久化方面形成系统化突破，输出兼具交互自然性与执行推动力的解决方案。
    
    \par （3）在系统开发与验证方面，各阶段节点需科学安排、具有可执行性。要求严格按照“需求分析-架构设计-模块开发-系统集成测试”的逻辑主线开展工作。特别是在后期的验证环节，需精心设计智能体工作流稳定性测试与端到端功能验证，确保多轮对话状态管理的可靠性与工具调用的准确性；需完成系统性能基准测试，验证高并发场景下的响应延迟与资源占用情况。最终需按时高质量交付包含系统设计说明书、可部署运行的完整系统、以及系统测试报告在内的全套成果。
}

\mbox{} \vfill

\signature{指导教师（签名）}
% comment the line above and uncomment the line below if you want to set a signature with a specific date.
% \signaturewithdate{指导教师（签名）}{1897}{5}{21}
